\subsection{对称张量与 Banach 空间}

13.1.1. 设 E 是交换环 A 上的模。若 n 是满足 \(\geqslant0\) 的整数，则以 \(\mathsf{TS}^{n}(\mathbf{E})\) 表示 E 的 n 次对称张量模（A, III, p. 71 和 A, IV, § 5, no. 2）；各 \(\mathsf{TS}^{n}(\mathbf{E})\) 的直和记作 \(\mathsf{TS}(\mathbf{E})\)。若 r 是整数或符号 \(\infty\)、\(\omega\) 之一，则置

\[
\mathsf{TS}^{(r)}(\mathrm{E})=\bigoplus_{n\leqslant r}\mathsf{TS}^{n}(\mathrm{E})\quad\text{且}\quad\mathsf{TS}^{(r)+}(\mathrm{E})=\bigoplus_{1\leqslant n\leqslant r}\mathsf{TS}^{n}(\mathrm{E});
\]

因此，当 \(r \geqslant 0\) 时，\(\mathsf{TS}^{(r)}(\mathbf{E})\) 是 \(\mathsf{TS}^{(0)}(\mathbf{E}) = \mathbf{A}\) 与 \(\mathsf{TS}^{(r)+}(\mathbf{E})\) 的直和。有 \(\mathsf{TS}^{(1)+}(\mathbf{E}) = \mathbf{E}\) 以及 \(\mathsf{TS}^{(\infty)}(\mathbf{E}) = \mathsf{TS}^{(\omega)}(\mathbf{E}) = \mathsf{TS}(\mathbf{E})\)。

若 n 是满足 \(\geqslant0\) 的整数且 \(x\in E\)，则将 \(\gamma_{n}(x)\) 记作元素 \(x\otimes\cdots\otimes x\)，它属于 \(\mathsf{TS}^{n}(\mathbf{E})\)；当 \(n=0\) 时，\(\gamma_{0}(x)=1\)。

13.1.2（“多项式映射”）。假设 A 是无限整环且 E 是自由 A-模。令 F 为 A-模。映射 \(f\colon\mathbf{E}\to\mathbf{F}\) 若满足以下等价条件，则称为次数 \(n\) 的齐次多项式（参见 A, IV, § 5, no. 9）：

\begin{enumerate}
\def\labelenumi{\alph{enumi})}
\item
  存在一个线性映射 \(\tilde{f}\colon\mathsf{TS}^{n}(\mathbf{E})\to\mathbf{F}\)，使得 \(f(x)=\tilde{f}(\gamma_{n}(x))\) 对所有 \(x\in\mathbf{E}\) 成立。
\item
  存在一个多线性映射 \(u: \mathbf{E}^{n} \to \mathbf{F}\)，使得 \(f(x) = u(x, \ldots, x)\) 对所有 \(x \in \mathbf{E}\) 成立。
\end{enumerate}

若此条件成立，映射 \(\tilde{f}\) 满足 \(a)\)，因而唯一；若 \(t \in \mathsf{TS}^n(\mathbf{E})\)，则将 \(\langle f, t \rangle\) 记作元素 \(\tilde{f}(t)\)。若 \(u: \mathbf{E}^n \to \mathbf{F}\) 满足 \(b)\)，则相应的从 \(\otimes^n\mathbf{E}\) 到 \(\mathbf{F}\) 的线性映射与 \(\tilde{f}\) 在子模 \(\mathsf{TS}^n(\mathbf{E})\) 上重合。当 \(n!\) 在 \(\mathbf{A}\) 中可逆时，可以唯一地选取 u 为对称映射。

13.1.3. 设 E 是 K 上的 Banach 空间，F 是 K 上的分离多范数空间。令 U 为 E 中 0 的开邻域，\(f\colon\mathbf{U}\to\mathbf{F}\) 是类 \(\mathbf{C}^{r}\) 的映射（\(r\in\mathbf{N}_{\mathbb{K}}\)），且 \(t\) 是 \(\mathsf{TS}^{(r)}(\mathbf{E})\) 的元素，参见 13.1.1（应用于环 A=K）。令 \(k\) 为满足 \(\leqslant r\) 的整数，使得 \(t\in\mathsf{TS}^{(k)}(\mathbf{E})\)。可以唯一地将 \(t\) 和 \(f\) 分解为

\[
t=t_{0}+\dots+t_{k},\quad\text{其中}\quad t_{i}\in\mathrm{TS}^{i}(\mathrm{E})
\]

以及

\[
f = f _ {0} + \dots + f _ {k} + h,
\]

其中 \(f_i\) 是次数为 \(i\) 的连续齐次多项式（\(0 \leqslant i \leqslant k\)），而 \(h\) 在点 0 处与 0 有阶数 \(\geqslant k\) 的接触。于是置

\[
\langle f, t \rangle = \sum_ {i = 0} ^ {k} \langle f _ {i}, t _ {i} \rangle ;
\]

这是 F 中的元素，且与 k 的选取无关。

13.1.4. 设 E 和 E' 是两个 Banach 空间，U 是 E 中 0 的开邻域，且 \(\varphi\colon\mathbf{U}\to\mathbf{E}'\) 是类 \(\mathbf{C}^{r}\) 的映射，满足 \(\varphi(0)=0\)。令 \(t\in\mathsf{TS}^{(r)}(\mathbf{E})\)。存在唯一一个 \(t'\in\mathsf{TS}^{(r)}(\mathbf{E}')\)，使得对于任意分离多范数空间 F、任意开邻域 \(\mathbf{U}'\)（位于 \(\mathbf{E}'\) 中且含 0），以及任意 \(f\colon\mathbf{U}'\to\mathbf{F}\) 类 \(\mathbf{C}^{r}\) 的映射，都有

\[
\langle f, t ^ {\prime} \rangle = \langle f \circ \varphi , t \rangle .\tag{*}
\]

这个元素 \(t'\) 记作 \(\varphi_{*}(t)\)。映射

\[
\varphi_ {*} \colon \mathsf {T S} ^ {(r)} (\mathrm{E}) \to \mathsf {T S} ^ {(r)} (\mathrm{E} ^ {\prime})
\]

从而所定义的映射是线性的。
% semantic-fragments: legacy-input-seam
\relax{}
